%%%%%%%%%%%%%%%%%%%%%%%%%%%%%%%%%
%
%       EXERCÍCIO
%
%%%%%%%%%%%%%%%%%%%%%%%%%%%%%%%%%

\definevalorquestao

\begin{exercicioBanco}[\valorquestao]
Sejam \(A\) e \(B\) eventos de um espaço amostral, onde \(P(B) > 0\). Mostre que:
%
\begin{enumerate}[(a)]
    \item Se \(P(A \mid B) = P(A)\), então \(P(A \cap B) = P(A)P(B)\);
    %
    \item \(P(A^{\complement} \mid B) = 1 - P(A \mid B)\).
\end{enumerate}
%
\end{exercicioBanco}

